\documentclass[11pt,a4paper]{amsart}
\usepackage[a4paper,landscape,margin=20mm]{geometry}
\usepackage{amsmath,amssymb,amsthm,graphicx,xcolor}
\usepackage[no-math]{fontspec}
\setmainfont{lmroman10-regular.otf}[BoldFont=lmroman10-bold.otf,ItalicFont=lmroman10-italic.otf,BoldItalicFont=lmroman10-bolditalic.otf,SmallCapsFont=lmromancaps10-regular.otf]
\usepackage[inline]{enumitem}
\usepackage[utf8]{inputenc} % required by tex-lint
\usepackage{graphicx}
\usepackage{subcaption}
\usepackage{booktabs} % for tables without vertical rules

\usepackage[colorlinks=true,citecolor=blue,linkcolor=blue]{hyperref}
\usepackage{cleveref}

\usepackage{amsmath,amsthm,amssymb,amsfonts}
\usepackage{caption}
\usepackage{bbm}
\usepackage{multirow}
\usepackage{xspace}
\usepackage{float}

\usepackage{tikz}
\usetikzlibrary{arrows,shapes,automata, calc, chains, matrix, trees, positioning, scopes}
\usetikzlibrary{decorations.pathmorphing}
\tikzset{snake it/.style={decorate, decoration=snake}}


\newcommand{\OO}{\mathcal{O}}
\newcommand{\Aa}{\mathcal{A}}
\newcommand{\BB}{\mathcal{B}}
\newcommand{\CC}{\mathcal{C}}
\newcommand{\DD}{\mathcal{D}}
\newcommand{\WW}{\mathcal{W}}
\newcommand{\MM}{\mathcal{M}}
\newcommand{\ee}{\mathbf{e}}
\newcommand{\PP}{\mathcal{P}}
\newcommand{\RR}{\mathcal{R}}

\newcommand{\VV}{\mathcal{V}}
\newcommand{\Start}{\mathcal{S}}

\newcommand{\ACZ}{\textnormal{\textsf{AC$^0$}}\xspace}
\newcommand{\NCT}{\textnormal{\textsf{NC$^2$}}\xspace}
\newcommand{\NCThree}{\textnormal{\textsf{NC$^3$}}\xspace}
\newcommand{\NP}{\textnormal{\textsf{NP}}\xspace}
\newcommand{\NC}{\textnormal{\textsf{NC}}\xspace}
\newcommand{\coNP}{\textnormal{\textsf{coNP}}\xspace}
\newcommand{\PSPACE}{\textnormal{\textsf{PSPACE}}\xspace}
\newcommand{\NPSPACE}{\textnormal{\textsf{NPSPACE}}\xspace}
\newcommand{\NL}
{\textnormal{\textsf{NL}}\xspace}
\newcommand{\coNL}
{\textnormal{\textsf{coNL}}\xspace}
\newcommand{\LL}{\textnormal{\textsf{L}}\xspace}
\newcommand{\PTIME}{\textnormal{\textsf{P}}\xspace}

\newcommand{\KK}{\mathbb{K}}
\newcommand{\ZZ}{\mathbb{Z}}
\newcommand{\NN}{\mathbb{N}}

\DeclareMathOperator{\modu}{mod}


%\DeclareMathOperator{\rk}{rank_$\mathbb{B}$}
\newcommand{\rk}{\mathsf{rank}_{\mathbb{B}}}
\newcommand{\rkr}{\mathsf{rank}_{\mathbb{R}}}
\newcommand{\rku}{\mathsf{rank}_{\mathsf{un}}}
\newcommand{\rkg}{\mathsf{rank}}

\newcommand{\B}{\mathbb{B}}
\newcommand{\R}{\mathbb{R}}
\newcommand{\RNN}{\mathbb{R}_{\ge 0}}
\newcommand{\RP}{\mathbb{R}_{> 0}}


\newcommand{\spn}[1]{\mathord{\langle #1 \rangle}}
\newcommand{\Mer}{\mathsf{Mer}}
%\newcommand{\rk}{\mathit{rk}}
\newcommand{\wt}[1]{\mathsf{wt}(#1)}
\newcommand{\mcw}{\mathsf{mcw}}
\newcommand{\mrw}{\mathsf{mrw}}
\newcommand{\MCol}{\mathsf{MCol}}
\newcommand{\MRow}{\mathsf{MRow}}
\newcommand{\MatMult}{\mathsf{MatMult}}


\newcommand{\new}[1]{\textcolor{red}{#1}}

\tikzset{AUT style/.style={>=angle 60,every edge/.append style={},every state/.style={minimum size=16,inner sep=0}}}


% the following environments switch to a slanted font:
\theoremstyle{plain}

\newtheorem{thm}{Theorem}[section]
\newtheorem{cor}[thm]{Corollary}
\newtheorem{lem}[thm]{Lemma}
\newtheorem{slem}[thm]{Sublemma}
\newtheorem{prop}[thm]{Proposition}

% the following environments keep the roman font:
\theoremstyle{definition}

\newtheorem{asm}[thm]{Assumption}
\newtheorem{rem}[thm]{Remark}
\newtheorem{rems}[thm]{Remarks}
\newtheorem{exa}[thm]{Example}
\newtheorem{exas}[thm]{Examples}
\newtheorem{defi}[thm]{Definition}
\newtheorem{conv}[thm]{Convention}
\newtheorem{conj}[thm]{Conjecture}
\newtheorem{prob}[thm]{Problem}
\newtheorem{oprob}[thm]{Open Problem}
\newtheorem{oprobs}[thm]{Open Problems}
\newtheorem{algo}[thm]{Algorithm}
\newtheorem{obs}[thm]{Observation}
\newtheorem{desc}[thm]{Description}
\newtheorem{fact}[thm]{Fact}
\newtheorem{qu}[thm]{Question}
\newtheorem{oqu}[thm]{Open Question}
\newtheorem{pty}[thm]{Property}
\newtheorem{clm}[thm]{Claim}
\newtheorem{nota}[thm]{Notation}
\newtheorem{com}[thm]{Comment}
\newtheorem{coms}[thm]{Comments}

%% goal is to omit parens from note, place cite as part of theorem number
\newtheoremstyle{defC}%
  {6pt}% (space above)
  {6pt}% (space below)
  {\normalfont}% (body font)
  {}% (indent amount)
  {\bfseries}% {theorem head font}
  {{\bfseries .}}% {punctuation after theorem head}
  {5pt plus 1pt minus 1pt}% {space after theorem head}
  {\thmname{#1} \thmnumber{#2} \thmnote{\normalfont#3}}% {theorem head spec}
\theoremstyle{defC}
\newtheorem{defiC}[thm]{Definition}
\newtheorem{remC}[thm]{Remark}
\newtheorem{exaC}[thm]{Example}

\newtheoremstyle{thmC}%
  {6pt}% (space above)
  {6pt}% (space below)
  {\itshape}% (body font)
  {}% (indent amount)
  {\bfseries}% {theorem head font}
  {{\bfseries .}}% {punctuation after theorem head}
  {5pt plus 1pt minus 1pt}% {space after theorem head}
  {\thmname{#1} \thmnumber{#2} \thmnote{\normalfont#3}}% {theorem head spec}

\theoremstyle{thmC}
\newtheorem{thmC}[thm]{Theorem}
\newtheorem{propC}[thm]{Proposition}
\newtheorem{lemC}[thm]{Lemma}

% switch back to standard theorem style for document specific theorems
\theoremstyle{plain}

\numberwithin{equation}{section}
\AtBeginDocument{\def\shorttitle{The rank of unambiguous finite automata}}
